\chapter{Production-Readiness and Governance Checklists}
\index{production readiness}\index{checklists}\index{governance!checklists}%
This appendix distills the production notes, pitfalls, and acceptance criteria scattered through
the chapters into domain checklists. Use them two ways: as a pre-launch gate (every box ticked or
consciously deferred with a reason) and as exam preparation---Professional Data Engineer scenario
questions are, in large part, these checklists wearing a story. Items marked \emph{(neg)} require
a negative test: proving the control denies what it must deny.

\section{Storage and Data Lake (Chapters 1, 8)}
\begin{itemize}
    \item[$\square$] Bucket location matches compute locality; dual/multi-region only where an
    availability requirement justifies the price.
    \item[$\square$] Lifecycle rules transition and expire objects; no bucket grows unbounded
    (Archive minimum-duration cost considered before archiving active data).
    \item[$\square$] Versioning, retention policy, or soft delete chosen deliberately per bucket;
    the three mechanisms are not conflated.
    \item[$\square$] Uniform bucket-level access enabled; no legacy ACLs leaking object-level
    exceptions.
    \item[$\square$] Encryption posture recorded: Google-managed, CMEK, or CSEK, with key rotation
    and \emph{(neg)} access revocation tested.
    \item[$\square$] Transfer jobs (Storage Transfer Service) validated on bytes, object counts,
    sampled checksums, and incremental sync---not just ``it ran.''
    \item[$\square$] Lake zones mapped to a medallion layout (\texttt{raw/}, \texttt{validated/},
    \texttt{curated/}) with IAM boundaries per zone.
    \item[$\square$] BigLake connections enforce policy tags and row access policies on lake files;
    \emph{(neg)} a direct GCS read by an unauthorized identity fails.
\end{itemize}

\section{Warehousing and Analytics (Chapters 5--8)}
\begin{itemize}
    \item[$\square$] Fact tables declare a grain; every measure and key respects it.
    \item[$\square$] Large tables partitioned (usually by date) and clustered on the dominant
    high-cardinality filter columns; bytes-billed before/after measured.
    \item[$\square$] Materialized views cover the top recurring aggregations; refresh cost is
    budgeted.
    \item[$\square$] Pricing model chosen from measured slot usage
    (\texttt{INFORMATION\_SCHEMA.JOBS\_BY\_PROJECT}), not guessed: on-demand for spiky, Editions
    reservation for steady workloads.
    \item[$\square$] Custom query cost quotas per user/project; dry-run byte checks in CI.
    \item[$\square$] SCD Type~2 dimensions match on key \emph{and} \texttt{is\_current}; history
    reconstruction tested on a known entity.
    \item[$\square$] External/BigLake tables accepted with their limits (no clustering, no managed
    statistics) or promoted to native tables where performance matters.
    \item[$\square$] Semantic layer (LookML or equivalent) owns metric definitions; no dashboard
    re-implements a governed measure.
    \item[$\square$] BI Engine capacity sized for hot dashboard tables only, not batch ELT.
    \item[$\square$] Analytics Hub listings grant read-only linked datasets; revocation tested.
\end{itemize}

\section{Batch ETL and Orchestration (Chapters 9--12, 24)}
\begin{itemize}
    \item[$\square$] Dataproc clusters are ephemeral per job (or serverless batches), with
    preemptible capacity confined to secondary workers and jobs made retry-safe.
    \item[$\square$] Spark jobs sized from evidence: skew inspected in the Spark UI before memory
    is raised; autoscaling policy bounds and cooldowns set.
    \item[$\square$] Data Fusion instances use private IP where sources are internal; macros
    parameterize dates and paths so one definition serves schedules, backfills, and reruns.
    \item[$\square$] Datastream streams monitored for freshness lag; backfill strategy documented
    per source table.
    \item[$\square$] Every Airflow task is idempotent (deterministic partition overwrites or
    upserts); retries and clears are safe by construction.
    \item[$\square$] XComs carry small control values only---paths, counts, verdicts---never
    payloads.
    \item[$\square$] Event-driven readiness uses sensors (file arrival) rather than clock guesses;
    schedules have catchup behavior chosen deliberately.
    \item[$\square$] Dataform (or dbt) assertions gate the pipeline; a failing test blocks
    dependents.
    \item[$\square$] Terraform state is remote and locked; applies run from CI after review;
    Workload Identity Federation replaces all committed keys.
    \item[$\square$] Schema evolution follows expand--contract: additive change, dual-run shim,
    consumer migration, contract cleanup---never a one-release breaking change.
\end{itemize}

\section{Streaming and Real-Time (Chapters 13--16)}
\begin{itemize}
    \item[$\square$] Topics carry Avro/Proto schemas so the broker rejects contract violations at
    publish time.
    \item[$\square$] Every subscription has a dead-letter topic and a triage owner; poison
    messages never silently vanish.
    \item[$\square$] Deduplication is designed in: message IDs, idempotent sinks, or Dataflow
    exactly-once mode---at-least-once delivery is assumed.
    \item[$\square$] Pub/Sub Lite chosen only after the steady-throughput cost comparison against
    Standard.
    \item[$\square$] \texttt{oldest\_unacked\_message\_age} alerts exist per subscription.
    \item[$\square$] Windowing, watermark, and allowed-lateness choices are written down per
    pipeline; late-data handling (side output or update) is explicit.
    \item[$\square$] Per-key state is bounded by timers or watermark garbage collection; unbounded
    state growth is treated as a bug.
    \item[$\square$] Sinks are exactly-once (Storage Write API) or at-least-once with a merge/dedup
    step---and the choice is justified by cost.
    \item[$\square$] Streaming jobs drain (not cancel) for deploys; autoscaling floor/ceiling set
    for known traffic peaks.
    \item[$\square$] Reverse-ETL syncs resolve identity before writing to operational systems and
    respect API rate limits on a freshness-versus-cost cadence.
\end{itemize}

\section{Machine Learning and MLOps (Chapters 17--20)}
\begin{itemize}
    \item[$\square$] The simplest sufficient model is the baseline (BQML logistic regression or
    AutoML) before any custom training spend.
    \item[$\square$] Training/serving skew prevented by construction: one feature definition
    (feature store or shared SQL) feeds both paths.
    \item[$\square$] Evaluation uses the metric that matches the decision---precision/recall at a
    chosen threshold for imbalanced problems, never bare accuracy.
    \item[$\square$] Models live in a registry with lineage to dataset and code; promotion is an
    alias move; rollback is tested.
    \item[$\square$] Online endpoints serve only latency-bound paths; batch prediction covers bulk
    scoring; endpoints are undeployed when idle.
    \item[$\square$] Retraining triggers fire on drift evidence (with data-quality preconditions),
    not on a calendar; candidates must beat the incumbent before registration.
    \item[$\square$] External AI API calls inside pipelines are batched, bounded in concurrency,
    retried with backoff and jitter, and dead-lettered on permanent failure; per-call cost is
    forecast and alerted.
    \item[$\square$] GenAI/RAG pipelines govern source chunks and prompts like data assets:
    versioned, access-controlled, evaluated.
\end{itemize}

\section{Security, Governance, and Compliance (Chapters 21--23)}
\begin{itemize}
    \item[$\square$] Least privilege everywhere: predefined or custom roles, no primitive Owner/
    Editor grants in shared projects; service accounts per workload, not shared.
    \item[$\square$] No service-account key files exist; Workload Identity (Federation) issues
    short-lived tokens instead.
    \item[$\square$] VPC Service Controls perimeters enclose sensitive projects; access levels
    carve only documented exceptions; \emph{(neg)} a valid-credential request from outside the
    perimeter is denied and alerted.
    \item[$\square$] CMEK on regulated datasets and buckets with rotation; key destruction
    procedure rehearsed (it is the erasure backstop).
    \item[$\square$] Cloud Audit Logs (including Data Access where required) route to a log
    bucket or BigQuery with alerting on policy denials and IAM changes.
    \item[$\square$] Dataplex lakes/zones map to data-mesh domains; domains own product quality
    and SLAs; the platform owns taxonomy and minimum rules.
    \item[$\square$] Policy tags (not per-role views) enforce column-level security; row access
    policies have \emph{(neg)} tests.
    \item[$\square$] AutoDQ scans plus declared quality rules cover completeness, uniqueness,
    validity, and freshness on curated assets; findings route to owners.
    \item[$\square$] DLP discovery scans the estate on a schedule; sensitive columns are masked
    (irreversible) or tokenized (KMS-wrapped, reversible under audit)---never left cleartext
    ``temporarily.''
    \item[$\square$] Right-to-erasure design accounts for time travel, backups, and replicas;
    quasi-identifier risk analysis runs on ``safe'' de-identified tables.
    \item[$\square$] Lineage is queryable from dashboard tile back to source system; impact
    analysis is a lookup, not archaeology.
\end{itemize}

\section{Reliability and FinOps (Chapter 23 and throughout)}
\begin{itemize}
    \item[$\square$] RPO/RTO stated per pipeline; restore and failover are drilled, not assumed.
    \item[$\square$] Budgets alert at 50/90/100\%; billing export to BigQuery feeds a weekly
    cost-by-service review.
    \item[$\square$] Idle-billing resources (clusters, Composer, endpoints, instances) have
    same-day teardown or auto-stop.
    \item[$\square$] Every architecture review answers: what breaks when this component fails, who
    is paged, and what does it cost at 10$\times$ volume.
\end{itemize}

\begin{tipbox}
During the exam, each checklist above compresses into one elimination rule: prefer the managed
service, the least-privilege grant, the answer that includes monitoring, and the design that
satisfies the stem's stated constraint---then verify none of these checklists is violated.
\end{tipbox}
